\section{Introdução}

A crescente digitalização de serviços e a onipresença de dispositivos computacionais pessoais gera um volume de dados distribuídos sem precedentes,
o que impulsiona o desenvolvimento de soluções baseadas em aprendizado de máquina (ML) \cite{li2022federated, kim2023kfl}.
Em cenários convencionais, uma solução de ML exigiria a centralização dessas informações em um único servidor para o treinamento.
No entanto, a típica natureza sensível dos dados distribuídos, como é o caso na área da saúde,
impõe restrições de privacidade que impedem essa abordagem centralizada \cite{mcmahan2017communication, zhang2025lcfed}.

Para superar esse obstáculo, o Aprendizado Federado (do inglês, \textit{Federated Learning} -- FL)
surgiu como um paradigma onde cada cliente recebe uma cópia de um modelo global inicial,
realiza o treinamento com seus próprios dados e transmite apenas as atualizações do modelo \cite{mcmahan2017communication}.
O servidor, por sua vez, agrega essas atualizações com uma média ponderada (FedAvg) para criar uma nova versão do modelo global.
Esse processo se repete iterativamente por um determinado número de vezes até a convergência do modelo global \cite{mcmahan2017communication, li2022federated}.

Apesar dos avanços proporcionados pelo FL, sua eficácia é frequentemente desafiada pela heterogeneidade estatística dos dados nos clientes (dados não-IID).
Essa não uniformidade impacta negativamente a agregação,
tornando o modelo global subótimo para a maioria dos participantes \cite{ghosh2020efficient, briggs2020federated}.

Para mitigar esse impacto, uma linha de pesquisa proeminente foca no Aprendizado Federado Clusterizado
(do inglês, \textit{Clustered FederatedLearning} -- CFL).
A intuição por trás do CFL é agrupar clientes com distribuições de dados semelhantes
e treinar modelos especializados para cada grupo \cite{ghosh2020efficient, briggs2020federated, zhang2025lcfed}.
Essa abordagem tem demonstrado melhorar significativamente o desempenho em ambientes altamente heterogêneos \cite{li2022federated, kim2023kfl}.

Briggs et al. \cite{briggs2020federated} introduziram o FL+HC,
que aplica agrupamento hierárquico sobre a distância euclidiana dos pesos das atualizações.
A principal limitação dessa estratégia reside no fato de que,
uma vez que o modelo é cindido em clusters independentes,
a troca de conhecimento entre eles é permanentemente interrompida.
Essa segmentação estrita impede que clientes em clusters distintos,
mas com características latentes comuns,
continuem a se beneficiar de uma base de conhecimento compartilhada.

Buscando otimizar o momento da clusterização,
o método K-FL \cite{kim2023kfl} emprega um Filtro de Kalman para monitorar a tendência da acurácia e decidir o instante ideal para o agrupamento.
Apesar da inovação no uso de filtros estocásticos para lidar com o ruído das atualizações,
o K-FL ainda mantém a premissa de fronteiras rígidas (\textit{hard boundaries}),
onde um dispositivo é sumariamente excluído de um cluster caso sua performance divirja da média,
uma decisão que pode ser prematura em cenários de alta variância nos dados.

Para mitigar o \textit{overhead} de comunicação,
Zhang et al. \cite{zhang2025lcfed} propuseram o LCFed,
que utiliza a divisão do modelo (\textit{model splitting}) e aproximações de baixo posto (\textit{low-rank}) para calcular similaridades de forma eficiente.
No entanto, o LCFed ainda restringe a decisão final do cliente a um único submodelo,
limitando a expressividade da personalização em casos onde o perfil de dados do cliente transita entre diferentes distribuições.

No algoritmo IFCA, proposto por Ghosh et al. \cite{ghosh2020efficient},
o servidor mantém múltiplos modelos globais e os transmite aos clientes.
Cada participante avalia os modelos recebidos e seleciona aquele que minimiza sua perda empírica local para realizar o treinamento.
Essa abordagem exige que o cliente se vincule estritamente a um único modelo por rodada,
o que pode negligenciar conhecimento potencialmente útil dos modelos de outros clusters.

Finalmente, o Aprendizado Federado com Agrupamento Suave (FLSC),
apresentado por Li et al. \cite{li2022federated}, possibilitou que um cliente pertença a múltiplos clusters simultaneamente.
Graças a isso, o FLSC demonstrou capturar melhor a natureza dos dados reais~\cite{...}.
Contudo, o método obriga cada cliente a participar de exatamente $N$ clusters, sendo $N$ um hiperparâmetro.
Além disso, o FLSC não considera a possibilidade de que o peso de participação do cliente em cada um dos $N$ clusters possa ser diferente.

Para superar essas limitações, este trabalho propõe o FLASC (\textit{Federated Learning with Adaptive Soft Clustering}).
Diferente do FLSC, o FLASC permite que a quantidade de clusters integrados varie conforme a necessidade de cada cliente.
Além disso, ele propõe uma adaptação mais precisa e robusta à heterogeneidade dos dados ao
ponderar a contribuição de cada modelo de cluster proporcionalmente à sua afinidade com os dados locais.
O restante deste trabalho está organizado da seguinte forma:
a Seção \ref{sec:proposta} descreve detalhadamente o funcionamento do algoritmo FLASC;
a Seção \ref{sec:metodologia} apresenta a metodologia de avaliação e o ambiente experimental;
a Seção \ref{sec:resultados} expõe e discute os resultados obtidos;
e, finalmente, a Seção \ref{sec:conclusao} traz as considerações finais e propostas de trabalhos futuros.
